\newpage
\section{Medios con respuesta lineal}
\label{sec.2-medioslineales}

reemplazamos la \(\rho\) y \(\mathbf{j}\) de las cargas puntuales por campos continuos que describen acumulación de carga y corriente dentro de un medio continuo.

\subsection{Marco general}
\label{subsec.2-general}

definir polarización y magnetización

ecuaciones de maxwell macroscópicas

relaciones constitutivas generales

particularizar a medios isótropos, homogéneos, locales.


\subsection{Cristales fotónicos}
\label{subsec.2-cristales}

novotny's derivation for 1d

intuition behind what bloch modes are (joannopoulos, p. 40)

The Bloch state Hk (r)e-iomegat
is a plane wave ei(k·r-omegat) multiplied by a periodic “envelope” function uk (r). It
propagates through the crystal without scattering, because k is conserved (apart
from addition of reciprocal lattice vectors, which is merely a relabeling)

In the familiar case of a homogeneous, isotropic medium, k is the direction in
which the wave propagates, but this is not necessarily true in a periodic medium.
Rather, the direction and the speed with which electromagnetic energy passes
through the crystal are given by the group velocity v, which is a function of both
the band index n and the wave vector k:


We can index the eigenstates (psikn or Hkn) using the translation operator
eigenvalues. These can be labelled in terms of the wave vectors and bands in
the Brillouin zone. All of the eigenstates can be cast in Bloch form: a periodic
function modulated by a plane wave. The field can propagate through the crystal
in a coherent manner, as a Bloch wave. The understanding of Bloch waves for
electrons explained one of the great mysteries of nineteenth-century physics: Why
do electrons behave like free particles in many examples of conducting crystals?
In the same way, a photonic crystal provides a synthetic medium in which light
can propagate, but in ways quite different from propagation in a homogeneous
medium

The bands above and below the gap can be distinguished by where the energy
of their modes is concentrated: in the high-eps regions, or in the low-eps regions. Often,
especially in the two- and three-dimensional crystals of the later chapters, the
low-eps regions are air regions. For this reason, it is convenient to refer to the
band above a photonic band gap as the air band, and the band below a gap as
the dielectric band. The situation is analogous to the electronic band structure

The key observation of the section The Physical Origin of Photonic Band Gaps
was that the periodicity of the crystal induced a gap in its band structure. No
electromagnetic modes are allowed to have frequencies in the gap. But if this is
indeed the case, what happens when we send a light wave (with a frequency in
the photonic band gap) onto the face of the crystal from outside?

2d and 3d?

\subsection{Metamateriales}
\label{subsec.2-metamateriales}


\subsection{Guías de onda}
\label{subsec.2-guias}